\documentclass[12pt, a4paper]{article}

% ── Packages ──────────────────────────────────────────────────────────────────
\usepackage[utf8]{inputenc}
\usepackage[T1]{fontenc}
\usepackage{amsmath, amssymb, amsthm}
\usepackage{geometry}
\usepackage{hyperref}
\usepackage{booktabs}
\usepackage{enumitem}
\usepackage{algorithm}
\usepackage{algpseudocode}
\usepackage{graphicx}
\usepackage{xcolor}
\usepackage{caption}
\usepackage[backend=biber, style=authoryear]{biblatex}

\addbibresource{references.bib}

\geometry{margin=2.5cm}
\hypersetup{colorlinks=true, linkcolor=blue!60!black, citecolor=blue!60!black, urlcolor=blue!60!black}

\newtheorem{definition}{Definition}
\newtheorem{remark}{Remark}

% ── Title ─────────────────────────────────────────────────────────────────────
\title{Construction of the SPX Implied Volatility Surface\\[6pt]
       \large Step~1 of the Local Stochastic Volatility Model}
\author{Imperial College London --- Master's Thesis}
\date{\today}

\begin{document}
\maketitle
\tableofcontents
\newpage

% ══════════════════════════════════════════════════════════════════════════════
\section{Introduction and Motivation}
\label{sec:intro}
% ══════════════════════════════════════════════════════════════════════════════

The goal of this chapter is to construct a continuous implied volatility surface
from discrete market observations of S\&P~500 (SPX) index options. This surface
serves as the calibration target for the Local Stochastic Volatility (LSV) model
developed in later chapters.

An \emph{implied volatility surface} is a two-dimensional function
$\sigma_{\text{imp}}(K, T)$ that maps every combination of strike price~$K$ and
time to maturity~$T$ to a single volatility number. Rather than assuming that
volatility is constant (as the original Black--Scholes model does), the surface
captures how the market prices risk differently depending on how far an option's
strike is from the current price and how much time remains until expiry
\parencite{gatheral2006volatility, cont2002dynamics}.

The pipeline proceeds in seven stages:
\begin{enumerate}[nosep]
    \item Fetch market data: spot price, risk-free rate, dividend yield, and the
          full SPX option chain.
    \item Clean the data: remove illiquid contracts and those violating
          no-arbitrage bounds.
    \item Filter to out-of-the-money options only.
    \item Invert the Black--Scholes formula to recover implied volatilities.
    \item Interpolate the scattered implied volatilities onto a regular grid
          using radial basis functions.
    \item Plot the resulting surface and per-expiry volatility smiles.
    \item Validate the surface by repricing a sample of market options.
\end{enumerate}

% ══════════════════════════════════════════════════════════════════════════════
\section{Market Data}
\label{sec:data}
% ══════════════════════════════════════════════════════════════════════════════

All data is sourced from Yahoo Finance via the \texttt{yfinance} Python library.
No paid API key is required; however, the data reflects delayed or end-of-day
quotes rather than real-time feeds.

\subsection{Spot Price}

The spot price $S$ is the most recent closing value of the S\&P~500 index
(\texttt{\^{}SPX}). This single number anchors the entire surface: all strikes
are measured relative to~$S$ through the moneyness ratio $K/S$.

\subsection{Risk-Free Rate}

The risk-free rate $r$ is obtained from the 13-week US Treasury Bill yield
(\texttt{\^{}IRX} on Yahoo Finance). The T-Bill is the standard short-term
risk-free proxy in academic option pricing because US Treasury securities carry
negligible default risk.

Yahoo Finance reports the T-Bill yield as an annualised percentage (e.g.,
$4.35\%$). We convert to a decimal and treat it as a continuously compounded
rate:
\begin{equation}
    r = \frac{\text{IRX}}{100}.
    \label{eq:rate}
\end{equation}

\begin{remark}
A single short-term rate is used for all maturities. In principle, one could use
a term structure of rates (different rates for different option expiries), which
would reduce the call--put implied volatility discontinuity discussed in
Section~\ref{sec:otm}. This simplification is standard in academic
implementations and any systematic effect is small relative to other sources of
noise.
\end{remark}

\subsection{Dividend Yield}

The continuous dividend yield $q$ is Yahoo Finance's trailing twelve-month
dividend yield for SPX. It enters the Black--Scholes formula through the forward
price:
\begin{equation}
    F = S \, e^{(r - q) T},
    \label{eq:forward}
\end{equation}
which represents the cost-of-carry-adjusted expected level of the index at
maturity~$T$. The term $r - q$ reflects the net cost of holding the index:
earning the risk-free rate on cash but forgoing the dividends paid by the index.

If the fetch fails, a fallback value of $q = 1.195\%$ is used, based on recent
historical SPX dividend yields.

\subsection{Option Chain}

Yahoo Finance provides full option chains for SPX across all available expiry
dates. For each expiry, both call and put contracts are retrieved. Each contract
carries the following fields (among others):
\begin{itemize}[nosep]
    \item \textbf{Strike} ($K$): the price at which the option can be exercised.
    \item \textbf{Bid / Ask}: the best quoted prices to sell / buy the option.
    \item \textbf{Open Interest}: the number of outstanding contracts, a proxy
          for liquidity.
    \item \textbf{Implied Volatility} (\texttt{iv\_yf}): Yahoo's own IV
          estimate, retained for comparison but not used in our surface.
\end{itemize}

The \textbf{mid price} is computed as the average of the bid and ask:
\begin{equation}
    \text{mid} = \frac{\text{bid} + \text{ask}}{2}.
    \label{eq:mid}
\end{equation}
This is used as the market price for all subsequent calculations because neither
the bid nor the ask alone represents a fair value---the mid is the best available
estimate of where a trade would execute.

\textbf{Time to maturity} is calculated as:
\begin{equation}
    T = \frac{\text{expiry date} - \text{today}}{\text{365 days}},
    \label{eq:ttm}
\end{equation}
expressed in years. Only options with $T \in [7/365,\; 2.0]$ are retained.
Options expiring within one week are excluded because their prices are dominated
by microstructure effects (gamma risk, pin risk), while options beyond two years
are too illiquid to provide reliable implied volatilities.

\textbf{Moneyness} is defined as:
\begin{equation}
    m = \frac{K}{S}.
    \label{eq:moneyness}
\end{equation}
A moneyness of~$1$ means the strike equals the current index level
(at-the-money). Values below~$1$ correspond to strikes below the index (in-the-money
calls / out-of-the-money puts), and values above~$1$ correspond to strikes above
the index.

% ══════════════════════════════════════════════════════════════════════════════
\section{Data Cleaning}
\label{sec:cleaning}
% ══════════════════════════════════════════════════════════════════════════════

Raw option data from Yahoo Finance contains many contracts that are stale,
illiquid, or mispriced. Using them would introduce noise into the implied
volatility surface. Three sequential filters are applied.

\subsection{Liquidity Filter}

The following conditions must all hold for an option to be retained:
\begin{enumerate}[nosep]
    \item \textbf{Open interest} $\geq 100$. Contracts with fewer outstanding
          positions are likely stale or thinly traded, producing unreliable mid
          prices.
    \item \textbf{Bid} $> 0$. A zero bid means no market participant is willing
          to buy the option at any price, so the quoted ask is meaningless.
    \item \textbf{Relative bid--ask spread} $\leq 20\%$:
          \begin{equation}
              \frac{\text{ask} - \text{bid}}{\text{mid}} \leq 0.20.
              \label{eq:spread}
          \end{equation}
          Wide spreads indicate poor liquidity; the mid price of such options is
          an unreliable estimate of fair value.
    \item \textbf{Moneyness} $K/S \in [0.70,\, 1.30]$. Deep out-of-the-money
          options (far from the current price) have near-zero prices and
          negligible vega, making implied volatility extraction numerically
          unstable.
    \item \textbf{Mid price} $> 0$. Options with zero or negative mid prices
          are discarded.
\end{enumerate}

\subsection{No-Arbitrage Filter}
\label{sec:noarb}

European option prices must satisfy model-independent bounds that follow from
the absence of arbitrage. Any option violating these bounds contains a data
error or has been mispriced by the market maker. The bounds are (with a $5\%$
tolerance to accommodate market microstructure):

\textbf{Calls:}
\begin{align}
    C &\geq 0.95 \times \max\!\big(0,\; S\,e^{-qT} - K\,e^{-rT}\big),
    \label{eq:call_lb} \\
    C &\leq 1.05 \times S\,e^{-qT}.
    \label{eq:call_ub}
\end{align}

\textbf{Puts:}
\begin{align}
    P &\geq 0.95 \times \max\!\big(0,\; K\,e^{-rT} - S\,e^{-qT}\big),
    \label{eq:put_lb} \\
    P &\leq 1.05 \times K\,e^{-rT}.
    \label{eq:put_ub}
\end{align}

The lower bound is the discounted intrinsic value: the minimum the option must
be worth to prevent riskless profit. The upper bound reflects the maximum
possible payoff: a call can never be worth more than the (discounted) index
itself, and a put can never be worth more than the (discounted) strike.

The $5\%$ tolerance allows for small bid--ask-induced deviations that occur
naturally in real markets without indicating a true arbitrage.

\subsection{Out-of-the-Money Filter}
\label{sec:otm}

After the above filters, we retain only \textbf{out-of-the-money} (OTM) options:
\begin{itemize}[nosep]
    \item \textbf{Calls} with $K \geq S$ (moneyness $\geq 1$),
    \item \textbf{Puts} with $K < S$ (moneyness $< 1$).
\end{itemize}

OTM options are preferred for three reasons:
\begin{enumerate}[nosep]
    \item \textbf{Liquidity}: OTM options are more actively traded than their
          in-the-money (ITM) counterparts, producing tighter bid--ask spreads
          and more reliable mid prices.
    \item \textbf{Numerical stability}: ITM option prices are dominated by
          intrinsic value (the payoff if exercised immediately), with only a
          thin layer of time value on top. Extracting implied volatility from
          time value alone magnifies any noise in the mid price.
    \item \textbf{Convention}: This is the standard approach in both academic
          literature and industry practice. The CBOE VIX methodology
          uses OTM options exclusively \parencite{cboe2019vix}.
\end{enumerate}

By put--call parity (Section~\ref{sec:bsm}), a European call and put at the
same strike and expiry must have the same implied volatility. Therefore, using
OTM options from both sides of the strike spectrum (puts on the left, calls on
the right) produces a single coherent surface without redundancy.

\begin{remark}
Even with OTM filtering, a small discontinuity at the at-the-money boundary
($K = S$) can appear because the single risk-free rate and dividend yield used
in the Black--Scholes inversion may not perfectly match the market-implied
forward price. This effect grows with maturity because the forward price
$F = S\,e^{(r-q)T}$ amplifies any error in $(r - q)$ over longer horizons.
Using a market-implied forward (extracted from put--call parity across strikes)
would eliminate this artefact.
\end{remark}

% ══════════════════════════════════════════════════════════════════════════════
\section{The Black--Scholes--Merton Model}
\label{sec:bsm}
% ══════════════════════════════════════════════════════════════════════════════

The Black--Scholes--Merton (BSM) model, introduced by
\textcite{black1973pricing} and extended by \textcite{merton1973theory} to
include continuous dividends, provides a closed-form formula for the price of a
European option.

\subsection{Assumptions}

The model assumes:
\begin{enumerate}[nosep]
    \item The underlying asset price $S_t$ follows a \textbf{geometric Brownian
          motion}:
          \begin{equation}
              dS_t = (r - q)\,S_t\,dt + \sigma\,S_t\,dW_t,
              \label{eq:gbm}
          \end{equation}
          where $r$ is the risk-free rate, $q$ is the continuous dividend yield,
          $\sigma$ is the constant volatility, and $W_t$ is a standard Brownian
          motion (a continuous random walk).
    \item Markets are frictionless: no transaction costs, no taxes, continuous
          trading.
    \item The risk-free rate $r$ and dividend yield $q$ are constant.
    \item The option can only be exercised at maturity (European-style).
\end{enumerate}

Under these assumptions, the log-return of the asset over any interval is
normally distributed:
\begin{equation}
    \ln\!\frac{S_T}{S_0} \sim \mathcal{N}\!\left(\left(r - q - \tfrac{\sigma^2}{2}\right)T,\;\; \sigma^2 T\right).
    \label{eq:log_return}
\end{equation}

\subsection{The Pricing Formulae}

\begin{definition}[BSM European Option Price]
The price of a European call option with strike $K$ and maturity $T$ is:
\begin{equation}
    C = S\,e^{-qT}\,\Phi(d_1) - K\,e^{-rT}\,\Phi(d_2),
    \label{eq:bsm_call}
\end{equation}
and the price of a European put is:
\begin{equation}
    P = K\,e^{-rT}\,\Phi(-d_2) - S\,e^{-qT}\,\Phi(-d_1),
    \label{eq:bsm_put}
\end{equation}
where $\Phi(\cdot)$ is the cumulative distribution function (CDF) of the
standard normal distribution, and
\begin{align}
    d_1 &= \frac{\ln(S/K) + (r - q + \tfrac{1}{2}\sigma^2)\,T}{\sigma\sqrt{T}},
    \label{eq:d1} \\[4pt]
    d_2 &= d_1 - \sigma\sqrt{T}.
    \label{eq:d2}
\end{align}
\end{definition}

In words: $d_1$ and $d_2$ measure how many standard deviations the current
index level is from the strike, adjusted for drift and time. The terms
$S\,e^{-qT}$ and $K\,e^{-rT}$ are the \emph{present values} of the index and
strike respectively, discounted at the dividend yield and risk-free rate.

$\Phi(d_1)$ can be interpreted as the probability (under the asset measure)
that the option expires in the money, while $\Phi(d_2)$ is the same probability
under the risk-neutral measure.

\subsection{Put--Call Parity}

A fundamental relationship links European calls and puts at the same strike and
maturity:
\begin{equation}
    C - P = S\,e^{-qT} - K\,e^{-rT}.
    \label{eq:pcp}
\end{equation}
This holds regardless of the volatility model and implies that a call and put at
the same $(K, T)$ must yield the same implied volatility.

\subsection{Vega}

The sensitivity of the option price to volatility is called \textbf{vega}:
\begin{equation}
    \mathcal{V} = \frac{\partial V}{\partial \sigma}
                = S\,e^{-qT}\,\phi(d_1)\,\sqrt{T},
    \label{eq:vega}
\end{equation}
where $\phi(\cdot)$ is the standard normal probability density function (PDF)
and $V$ denotes either the call or put price (vega is identical for both, by
put--call parity).

Vega is always non-negative and is largest for at-the-money options with long
maturities. It approaches zero for deep out-of-the-money or very short-dated
options. This quantity is critical for the implied volatility inversion described
next: it serves as the derivative in Newton--Raphson iteration.

% ══════════════════════════════════════════════════════════════════════════════
\section{Implied Volatility Computation}
\label{sec:iv}
% ══════════════════════════════════════════════════════════════════════════════

\subsection{Definition}

\begin{definition}[Implied Volatility]
The implied volatility $\sigma_{\text{imp}}$ of an option with observed market
price $V^{\text{mkt}}$ is the unique value of $\sigma$ that, when substituted
into the BSM formula, reproduces the market price:
\begin{equation}
    V^{\text{BSM}}(S, K, T, r, q, \sigma_{\text{imp}}) = V^{\text{mkt}}.
    \label{eq:iv_def}
\end{equation}
\end{definition}

There is no closed-form solution for $\sigma_{\text{imp}}$---the BSM formula
cannot be algebraically inverted for $\sigma$ (see \cite{jaeckel2015lets} for a
detailed treatment of numerical IV computation). Instead, we must solve
equation~\eqref{eq:iv_def} numerically. This is a root-finding problem: find
$\sigma$ such that
\begin{equation}
    f(\sigma) \coloneqq V^{\text{BSM}}(\sigma) - V^{\text{mkt}} = 0.
    \label{eq:rootfinding}
\end{equation}

\subsection{Newton--Raphson Method}
\label{sec:newton}

Newton--Raphson is an iterative root-finding algorithm. Given an initial guess
$\sigma_0$, each iteration updates the estimate using the function value and its
derivative:
\begin{equation}
    \sigma_{n+1} = \sigma_n - \frac{f(\sigma_n)}{f'(\sigma_n)}
                 = \sigma_n - \frac{V^{\text{BSM}}(\sigma_n) - V^{\text{mkt}}}
                                   {\mathcal{V}(\sigma_n)},
    \label{eq:newton}
\end{equation}
where $\mathcal{V}(\sigma_n)$ is the BSM vega at the current estimate. The
algorithm terminates when $|f(\sigma_n)| < \epsilon$, where
$\epsilon = 10^{-6}$ is the convergence tolerance.

\textbf{Why Newton--Raphson works well here}: The BSM price is a smooth, monotonically
increasing function of $\sigma$ (higher volatility always means a higher option
price). This means there is exactly one root, and Newton--Raphson converges
\emph{quadratically}---the number of correct digits roughly doubles each
iteration. Starting from a warm-start of $\sigma_0 = 0.25$ (a typical equity
volatility), convergence usually occurs in 3--5 iterations
\parencite{brenner1988note}.

\textbf{When it fails}: Newton--Raphson can stall when vega is very small
(deep OTM or very short-dated options), causing the update step
$f(\sigma)/\mathcal{V}(\sigma)$ to become excessively large. To handle this,
the estimate is clipped to the admissible range $[10^{-6},\; 5.0]$ at each
step, and if convergence is not achieved within 100 iterations, the algorithm
falls back to Brent's method.

\subsection{Brent's Method (Fallback)}
\label{sec:brent}

Brent's method \parencite{brent1973algorithms} is a bracketed root-finding
algorithm that combines bisection, secant, and inverse quadratic interpolation.
Unlike Newton--Raphson, it requires a bracket $[a, b]$ such that $f(a)$ and
$f(b)$ have opposite signs (guaranteeing a root exists between them), but it
does not require the derivative.

We use the bracket $[\sigma_{\text{lo}},\; \sigma_{\text{hi}}] = [10^{-6},\; 5.0]$.
If $f(\sigma_{\text{lo}})$ and $f(\sigma_{\text{hi}})$ do not have opposite
signs, the option price is inconsistent with the BSM model (e.g., below
intrinsic value), and the implied volatility is recorded as undefined (\texttt{NaN}).

\subsection{Sanity Checks}

Before attempting inversion, the market price is checked against the discounted
intrinsic value:
\begin{equation}
    \text{intrinsic} = \max\!\Big(0,\;
        \begin{cases}
            S\,e^{-qT} - K\,e^{-rT} & \text{(call)}, \\
            K\,e^{-rT} - S\,e^{-qT} & \text{(put)}.
        \end{cases}\Big)
    \label{eq:intrinsic}
\end{equation}
If $V^{\text{mkt}} \leq \text{intrinsic} + \epsilon$, the option has no time
value to extract, and implied volatility is undefined.

After inversion, options with $\sigma_{\text{imp}} \notin [0.01,\; 3.00]$ (i.e.,
outside $1\%$--$300\%$ annualised) are discarded as numerical artefacts.

% ══════════════════════════════════════════════════════════════════════════════
\section{Surface Interpolation}
\label{sec:interpolation}
% ══════════════════════════════════════════════════════════════════════════════

After implied volatility computation, we have a set of scattered observations
$\{(x_i, T_i, \sigma_i)\}_{i=1}^{N}$, where $x_i = \ln(K_i / S)$ is the
log-moneyness and $T_i$ is the time to maturity. These observations lie at
irregular positions determined by whatever strikes and expiries the exchange
lists. The downstream LSV model requires a \emph{continuous, smooth} surface
defined on a regular grid.

\subsection{Log-Moneyness}
\label{sec:logm}

We parameterise the strike axis using \textbf{log-moneyness}:
\begin{equation}
    x = \ln\!\left(\frac{K}{S}\right).
    \label{eq:log_moneyness}
\end{equation}
\begin{itemize}[nosep]
    \item $x = 0$: at-the-money ($K = S$).
    \item $x < 0$: strike below the index (OTM puts).
    \item $x > 0$: strike above the index (OTM calls).
\end{itemize}

Log-moneyness is preferred over raw strike $K$ or moneyness $K/S$ because:
\begin{enumerate}[nosep]
    \item The volatility smile is approximately symmetric in log-moneyness,
          improving interpolation quality.
    \item The at-the-money point is centred at $x = 0$ by construction.
    \item Numerical conditioning of the interpolation is improved (the range of
          $x$ is small and centred, rather than spanning thousands of strike
          dollars).
\end{enumerate}

\subsection{Radial Basis Function (RBF) Interpolation}
\label{sec:rbf}

Given $N$ scattered data points $\{(\mathbf{p}_i, \sigma_i)\}$ where
$\mathbf{p}_i = (T_i, x_i) \in \mathbb{R}^2$, we seek a smooth function
$\hat{\sigma}(\mathbf{p})$ that approximates the data. A \textbf{radial basis
function} (RBF) interpolant \parencite{buhmann2003radial} takes the form:
\begin{equation}
    \hat{\sigma}(\mathbf{p}) = \sum_{i=1}^{N} w_i \, \varphi\!\big(\|\mathbf{p} - \mathbf{p}_i\|\big) + c_0 + c_1 T + c_2 x,
    \label{eq:rbf}
\end{equation}
where:
\begin{itemize}[nosep]
    \item $\|\mathbf{p} - \mathbf{p}_i\|$ is the Euclidean distance from the
          evaluation point $\mathbf{p}$ to the $i$-th data point $\mathbf{p}_i$.
    \item $\varphi(r)$ is the radial basis function (kernel), which depends only
          on the distance $r$.
    \item $w_i$ are the weights, determined by fitting to the data.
    \item $c_0 + c_1 T + c_2 x$ is a linear polynomial term that captures any
          global trend (tilt) in the surface.
\end{itemize}

The key idea is that each data point acts as a ``bump'' centred at its location,
and the surface is built by summing these bumps with appropriate weights. The
word ``radial'' means each basis function is symmetric around its centre---it
depends only on the distance, not the direction.

\subsection{Thin-Plate Spline Kernel}
\label{sec:tps}

The specific kernel used is the \textbf{thin-plate spline}:
\begin{equation}
    \varphi(r) = r^2 \ln(r), \qquad r > 0, \qquad \varphi(0) = 0.
    \label{eq:tps_kernel}
\end{equation}

The thin-plate spline was introduced by \textcite{duchon1977splines} and is
extensively analysed in \textcite{wahba1990spline}. The name comes from a
physical analogy: if you take a thin sheet of metal and force it to pass through
a set of points, the shape it naturally assumes (the one that minimises bending
energy) is exactly the thin-plate spline. Formally, among all surfaces $f$ that
interpolate the data points, the thin-plate spline minimises the \emph{bending
energy}:
\begin{equation}
    \int\!\!\int \left[\left(\frac{\partial^2 f}{\partial x^2}\right)^{\!2}
    + 2\left(\frac{\partial^2 f}{\partial x\,\partial y}\right)^{\!2}
    + \left(\frac{\partial^2 f}{\partial y^2}\right)^{\!2}\right] dx\,dy.
    \label{eq:bending}
\end{equation}

This is critical for our application because:
\begin{enumerate}[nosep]
    \item The resulting surface is \textbf{infinitely differentiable} in the
          interior, which is required by Dupire's equation
          (Section~\ref{sec:dupire}) for computing local volatility. Dupire's
          formula involves both $\partial\sigma/\partial T$ and
          $\partial^2\sigma/\partial K^2$, so the surface must be at least
          twice differentiable.
    \item The surface is \textbf{as smooth as possible} given the data, which
          suppresses artificial oscillations that would produce spurious local
          volatility spikes.
    \item Unlike polynomial fitting, it does not require choosing a polynomial
          degree or risking Runge's phenomenon (wild oscillations at the
          boundaries).
\end{enumerate}

\subsection{Smoothing (Regularisation)}
\label{sec:smoothing}

With exact interpolation ($\hat{\sigma}(\mathbf{p}_i) = \sigma_i$ for all~$i$),
the surface would pass through every data point. However, since mid prices are
noisy estimates of true option values (they are simply the average of bid and
ask), the ``exact'' implied volatilities themselves contain noise. Forcing the
surface through noisy points would create artificial bumps and wiggles.

Instead, a \textbf{smoothing parameter} $\lambda > 0$ is introduced. The
weights $\{w_i\}$ are found by minimising:
\begin{equation}
    \sum_{i=1}^{N} \big[\hat{\sigma}(\mathbf{p}_i) - \sigma_i\big]^2
    + \lambda \cdot (\text{bending energy}),
    \label{eq:smoothing}
\end{equation}
which trades off closeness to the data (first term) against smoothness of the
surface (second term). Larger $\lambda$ produces a smoother surface that deviates
more from individual data points; smaller $\lambda$ tracks the data more closely
at the cost of potential oscillations.

The value used is $\lambda = N \times 5 \times 10^{-4}$, where $N$ is the
number of data points. This was chosen empirically to produce a surface with
sub-1-vol-point mean absolute error while remaining smooth enough for stable
local volatility computation.

\subsection{Grid Construction}

The interpolated surface is evaluated on a regular $60 \times 50$ grid:
\begin{itemize}[nosep]
    \item \textbf{Log-moneyness axis}: 60 equally spaced points between the
          2nd and 98th percentiles of the observed log-moneyness values.
    \item \textbf{TTM axis}: 50 equally spaced points between the 2nd and 98th
          percentiles of the observed times to maturity.
\end{itemize}

The 2nd--98th percentile trimming avoids extrapolation into regions with no
supporting data, where the thin-plate spline can produce unreliable values.
Any grid values below $0.01$ or above $3.0$ (annualised) are clipped as
numerical artefacts.

% ══════════════════════════════════════════════════════════════════════════════
\section{Connection to Local Volatility (Dupire's Equation)}
\label{sec:dupire}
% ══════════════════════════════════════════════════════════════════════════════

The primary consumer of the IV surface is Dupire's equation
\parencite{dupire1994pricing}, which converts the implied volatility surface
into a \textbf{local volatility} surface $\sigma_{\text{loc}}(K, T)$. Without
stating the full derivation (which belongs in the LSV model chapter), Dupire's
formula involves:
\begin{equation}
    \sigma_{\text{loc}}^2(K, T) = \frac{
        \dfrac{\partial C}{\partial T} + (r - q)\,K\,\dfrac{\partial C}{\partial K} + q\,C
    }{
        \dfrac{1}{2}\,K^2\,\dfrac{\partial^2 C}{\partial K^2}
    },
    \label{eq:dupire}
\end{equation}
where $C(K, T)$ is the market (or model) call price surface. Since $C$ is
computed from the IV surface via BSM, the smoothness of $\sigma_{\text{imp}}$
directly determines the stability of $\sigma_{\text{loc}}$. A noisy IV surface
would produce negative values inside the square root of~\eqref{eq:dupire},
yielding undefined local volatilities.

This is why the thin-plate spline with smoothing is essential: it provides
the twice-differentiable, physically sensible surface that Dupire's equation
requires. \textcite{fengler2009arbitrage} demonstrates that careful smoothing of
the IV surface is critical for obtaining stable and arbitrage-free local
volatilities, and \textcite{glaser2012arbitrage} further analyse approximation
methods for call price surfaces that preserve no-arbitrage conditions.
\textcite{homescu2011implied} provides a comprehensive survey of IV surface
construction methodologies, including kernel-based and spline-based approaches
consistent with the method used here.

% ══════════════════════════════════════════════════════════════════════════════
\section{Validation}
\label{sec:validation}
% ══════════════════════════════════════════════════════════════════════════════

To assess the quality of the interpolated surface, a random sample of 60 options
is drawn from the cleaned dataset. For each sampled option with known strike
$K_j$, maturity $T_j$, and computed implied volatility $\sigma_j$:

\begin{enumerate}[nosep]
    \item The surface IV is evaluated at $(\ln(K_j/S),\; T_j)$ using bilinear
          interpolation on the regular grid.
    \item The \textbf{IV absolute error} is recorded:
          \begin{equation}
              e_j^{\text{IV}} = |\hat{\sigma}(\ln(K_j/S),\; T_j) - \sigma_j|.
              \label{eq:iv_err}
          \end{equation}
    \item The option is \textbf{repriced} using the surface IV via the BSM
          formula, and the percentage price error is computed:
          \begin{equation}
              e_j^{\text{price}} = \frac{V^{\text{BSM}}(\hat{\sigma}_j) - V_j^{\text{mkt}}}{V_j^{\text{mkt}}} \times 100\%.
              \label{eq:price_err}
          \end{equation}
\end{enumerate}

\subsection{Metrics}

Two sets of metrics are reported:

\textbf{Primary (IV-space)}: The mean absolute error (MAE), root mean squared
error (RMSE), and maximum error of $e_j^{\text{IV}}$. These measure how
faithfully the continuous surface reproduces the discrete market-observed
implied volatilities. This is the correct metric for surface quality because the
surface exists in IV-space.

\textbf{Secondary (price-space, liquid options only)}: Percentage price errors
are only computed for options with market price $\geq \$10$. For cheap deep-OTM
options, vega is near zero, meaning a tiny IV difference causes an enormous
percentage price change. This is a mathematical singularity inherent to the
BSM mapping, not a deficiency of the surface. Excluding cheap options prevents
this artefact from dominating the error statistics.

A well-calibrated surface typically achieves an IV MAE of $0.5$--$1.0$ vol
points ($0.005$--$0.010$ in decimal), confirming that the smoothing preserves
the market structure while filtering out microstructure noise
\parencite{gatheral2006volatility}.

% ══════════════════════════════════════════════════════════════════════════════
\section{Summary of Outputs}
\label{sec:outputs}
% ══════════════════════════════════════════════════════════════════════════════

The pipeline produces the following files:

\begin{table}[h]
\centering
\begin{tabular}{@{}lp{8.5cm}@{}}
\toprule
\textbf{File} & \textbf{Description} \\
\midrule
\texttt{data/spx\_iv\_data.csv}        & Cleaned option chain with computed IVs \\
\texttt{data/validation\_results.csv}  & Repricing comparison for 60 sampled options \\
\texttt{arrays/ttm\_grid.npy}          & 1D array of 50 TTM grid values \\
\texttt{arrays/log\_m\_grid.npy}       & 1D array of 60 log-moneyness grid values \\
\texttt{arrays/iv\_surface.npy}        & $60 \times 50$ interpolated IV surface \\
\texttt{plots/iv\_surface\_3d.png}     & 3D surface visualisation \\
\texttt{plots/iv\_smiles.png}          & Per-expiry volatility smile overlays \\
\texttt{plots/validation.png}          & Validation diagnostic plots \\
\bottomrule
\end{tabular}
\caption{Output files produced by the IV surface pipeline.}
\label{tab:outputs}
\end{table}

The three \texttt{.npy} arrays constitute the surface in machine-readable form
and are loaded directly by the LSV model in Step~2.

% ══════════════════════════════════════════════════════════════════════════════
\section{Pipeline Summary}
\label{sec:pipeline}
% ══════════════════════════════════════════════════════════════════════════════

\begin{algorithm}[H]
\caption{SPX Implied Volatility Surface Construction}
\label{alg:pipeline}
\begin{algorithmic}[1]
\State Fetch spot price $S$, risk-free rate $r$, dividend yield $q$ from Yahoo Finance
\State Fetch full SPX option chain across all available expiries
\State \textbf{Filter}: remove options with low open interest, wide spreads, or zero bids
\State \textbf{Filter}: remove options violating no-arbitrage bounds
\State \textbf{Filter}: retain only OTM options (calls with $K \geq S$, puts with $K < S$)
\State Compute implied volatility $\sigma_i$ for each remaining option via Newton--Raphson / Brent
\State Discard options with $\sigma_i \notin [0.01, 3.0]$
\State Fit RBF thin-plate spline interpolant to scattered $(T_i, \ln(K_i/S), \sigma_i)$ data
\State Evaluate interpolant on regular $60 \times 50$ grid
\State Save surface arrays and plot diagnostics
\State Validate by repricing 60 random options and reporting IV MAE
\end{algorithmic}
\end{algorithm}

% ══════════════════════════════════════════════════════════════════════════════
\printbibliography[heading=bibintoc, title={References}]
% ══════════════════════════════════════════════════════════════════════════════

\end{document}
